\documentclass[11pt, a4paper]{article}

% Gemeinsame Style-Definitionen (TEXINPUTS muss templates/ enthalten — siehe scripts/build_tex.py)
% bfw-style.tex — gemeinsame Style-Definitionen für alle Handouts/Aufgabenhefte
% Voraussetzung: Kompilation mit xelatex (wegen fontspec)

\usepackage[ngerman]{babel}
\usepackage{geometry}
\geometry{a4paper, top=2.2cm, bottom=2.2cm, left=2.3cm, right=2.3cm}

\usepackage{fontspec}
% Fallback-Kette: Source Sans Pro -> Calibri -> Segoe UI -> Latin Modern Sans
% (Latin Modern ist immer mit MiKTeX vorhanden)
\IfFontExistsTF{Source Sans Pro}{%
  \setmainfont{Source Sans Pro}[Ligatures=TeX]%
}{\IfFontExistsTF{Calibri}{%
  \setmainfont{Calibri}[Ligatures=TeX]%
}{\IfFontExistsTF{Segoe UI}{%
  \setmainfont{Segoe UI}[Ligatures=TeX]%
}{%
  \setmainfont{Latin Modern Sans}[Ligatures=TeX]%
}}}
\IfFontExistsTF{Source Code Pro}{%
  \setmonofont{Source Code Pro}[Scale=0.92]%
}{\IfFontExistsTF{Consolas}{%
  \setmonofont{Consolas}[Scale=0.92]%
}{%
  \setmonofont{Latin Modern Mono}[Scale=0.92]%
}}

\usepackage{xcolor}
% Farbpalette: hell, freundlich, modern
\definecolor{bfwPrimary}{HTML}{0EA5A4}    % Petrol/Türkis - Primär
\definecolor{bfwPrimaryDark}{HTML}{0F766E} % dunkler Petrol für Headings
\definecolor{bfwAccent}{HTML}{F59E0B}     % Amber - Akzent
\definecolor{bfwSuccess}{HTML}{10B981}    % Grün - Erfolg / Lösung
\definecolor{bfwWarn}{HTML}{F97316}       % Orange - Achtung
\definecolor{bfwInk}{HTML}{0F172A}        % fast Schwarz
\definecolor{bfwInkSoft}{HTML}{475569}    % gedämpftes Grau
\definecolor{bfwBgSoft}{HTML}{F8FAFC}     % off-white
\definecolor{bfwBgTeal}{HTML}{ECFEFF}     % helles Türkis
\definecolor{bfwBgAmber}{HTML}{FEF3C7}    % helles Gelb
\definecolor{bfwBgRose}{HTML}{FFE4E6}     % helles Rosé für Eselsbrücken

\usepackage[colorlinks=true, linkcolor=bfwPrimaryDark, urlcolor=bfwPrimaryDark]{hyperref}
\usepackage{enumitem}
\setlist[itemize]{leftmargin=*, itemsep=2pt, topsep=2pt}
\setlist[enumerate]{leftmargin=*, itemsep=2pt, topsep=2pt}

\usepackage[most]{tcolorbox}
\usepackage{booktabs}
\usepackage{tikz}
\usetikzlibrary{decorations.pathmorphing, positioning, shapes.geometric, arrows.meta}

% Merkkästchen — wichtige Definition / Kernpunkt
\newtcolorbox{merksatz}[1][]{%
  enhanced, breakable,
  colback=bfwBgTeal,
  colframe=bfwPrimary,
  boxrule=0.6pt,
  arc=4pt,
  left=10pt, right=10pt, top=8pt, bottom=8pt,
  fonttitle=\bfseries\color{bfwPrimaryDark},
  title={\faLightbulb\ Merksatz},
  #1
}

% Eselsbrücke — Mnemoniks
\newtcolorbox{eselsbruecke}[1][]{%
  enhanced, breakable,
  colback=bfwBgRose,
  colframe=bfwAccent,
  boxrule=0.6pt,
  arc=4pt,
  left=10pt, right=10pt, top=8pt, bottom=8pt,
  fonttitle=\bfseries\color{bfwWarn},
  title={Eselsbrücke},
  #1
}

% Beispiel-Box
\newtcolorbox{beispiel}[1][]{%
  enhanced, breakable,
  colback=bfwBgSoft,
  colframe=bfwInkSoft,
  boxrule=0.4pt,
  arc=3pt,
  left=10pt, right=10pt, top=8pt, bottom=8pt,
  fonttitle=\bfseries\color{bfwInk},
  title={Beispiel},
  #1
}

% Lösungs-Box (für Aufgabenhefte)
\newtcolorbox{loesung}[1][]{%
  enhanced, breakable,
  colback=bfwBgAmber!40,
  colframe=bfwSuccess,
  boxrule=0.6pt,
  arc=3pt,
  left=10pt, right=10pt, top=8pt, bottom=8pt,
  fonttitle=\bfseries\color{bfwSuccess!70!black},
  title={Lösung},
  #1
}

% Glossar-Eintrag
\newcommand{\glossareintrag}[2]{%
  \par\noindent\textbf{\color{bfwPrimaryDark}#1} \enspace #2\par\smallskip
}

% Section-Stil — etwas farbiger als Standard
\usepackage{titlesec}
\titleformat{\section}
  {\Large\bfseries\color{bfwPrimaryDark}}
  {\thesection}{0.6em}{}
\titleformat{\subsection}
  {\large\bfseries\color{bfwPrimary}}
  {\thesubsection}{0.5em}{}
\titleformat{\subsubsection}
  {\normalsize\bfseries\color{bfwInk}}
  {\thesubsubsection}{0.4em}{}

% TikZ-Stil "skizzenhaft" — etwas weicher als Rough.js, aber stabil
\tikzset{
  roughbox/.style = {
    draw=bfwPrimaryDark, thick, fill=bfwBgTeal,
    rounded corners=4pt, inner sep=6pt
  },
  roughaccent/.style = {
    draw=bfwAccent, thick, fill=bfwBgAmber,
    rounded corners=4pt, inner sep=6pt
  },
  roughline/.style = {
    draw=bfwInkSoft, thick, -{Stealth[length=2.5mm]}
  }
}

% Bullet-Symbole (bewusst kein FontAwesome — vermeidet Font-Installations-Probleme)
\newcommand{\faLightbulb}{$\bullet$}
\newcommand{\faBookmark}{$\bullet$}

% Header/Footer
\usepackage{fancyhdr}
\pagestyle{fancy}
\fancyhf{}
\renewcommand{\headrulewidth}{0pt}
\fancyfoot[L]{\small\color{bfwInkSoft}\thepage}
\fancyfoot[R]{\small\color{bfwInkSoft} BFW M-064 Beschaffung im Büromanagement}


\title{\vspace*{-1.5cm}\Huge\bfseries\color{bfwPrimaryDark}Tag~1: Der Beschaffungsprozess\\[0.4em]
  \large\color{bfwInkSoft}\textnormal{Vom betrieblichen Bedarf zur Bezahlung}}
\author{\textsf{Beschaffung im Büromanagement}\\
\textsf{\small\copyright\ Can Siebert\,\textperiodcentered\,2026}}
\date{11.05.2026}

\begin{document}
\maketitle
\thispagestyle{empty}

\vspace{1em}
\begin{merksatz}[title={\bfseries Worum geht es heute?}]
Die Beschaffung ist das Bindeglied zwischen Beschaffungsmarkt und Produktion. Sie sorgt dafür,
dass die richtigen Materialien in der richtigen Menge und Qualität, am richtigen Ort
zur richtigen Zeit zum wirtschaftlichen Preis verfügbar sind --- und das nachhaltig.
In diesem Handout erarbeiten wir die Grundlagen, die Sie als Kauffrau bzw.\ Kaufmann für
Büromanagement im Berufsalltag täglich brauchen --- und die zugleich das Fundament für
die IHK-Abschlussprüfung Teil~1 und Teil~2 bilden.
\end{merksatz}

\tableofcontents
\vspace{1em}
\hrule
\vspace{1em}

\section{Beschaffung als Teil der Wertschöpfungskette}

Jedes Unternehmen verarbeitet Eingänge zu Ausgängen: Es nimmt Material, Wissen, Energie
und Arbeitskraft auf und wandelt diese in verkaufsfertige Produkte oder Dienstleistungen
um. Dieser geordnete Ablauf heißt \emph{Wertschöpfungskette}. Innerhalb dieser Kette
nimmt die Beschaffung eine Schlüsselrolle ein, denn sie steht ganz am Anfang und liefert
den Stoff, aus dem alle nachgelagerten Bereiche schöpfen.

Stellen Sie sich ein mittelständisches Unternehmen vor, das Aktenordner herstellt:
Wenn die Beschaffung morgens kein Karton bestellt hat, steht am Nachmittag die
Produktion still. Wenn sie zu teuer einkauft, kann das Endprodukt am Markt nicht mehr
wettbewerbsfähig angeboten werden. Wenn sie auf den falschen Lieferanten setzt, leiden
Liefertermine und Qualität --- die Folge sind Reklamationen und unzufriedene Kunden.
Beschaffungsfehler wirken also nicht isoliert, sondern entlang der gesamten Kette.

\begin{center}
\begin{tikzpicture}[
  every node/.style={font=\small\sffamily, align=center}
]
  \node[roughbox, minimum width=2.4cm, minimum height=1.1cm] (b) {Beschaffung};
  \node[roughaccent, minimum width=2.4cm, minimum height=1.1cm, right=0.6cm of b] (l) {Lagerung};
  \node[roughaccent, minimum width=2.4cm, minimum height=1.1cm, right=0.6cm of l] (p) {Produktion};
  \node[roughaccent, minimum width=2.4cm, minimum height=1.1cm, right=0.6cm of p] (v) {Vertrieb};
  \node[roughbox, minimum width=2.4cm, minimum height=1.1cm, right=0.6cm of v] (s) {Service};
  \draw[roughline] (b) -- (l);
  \draw[roughline] (l) -- (p);
  \draw[roughline] (p) -- (v);
  \draw[roughline] (v) -- (s);
\end{tikzpicture}
\end{center}
\vspace{-0.5em}
\begin{center}\small\itshape\color{bfwInkSoft}Die Wertschöpfungskette eines Industrieunternehmens. Die Beschaffung steht ganz vorn.\end{center}

\subsection{Die Beschaffung hat zwei Ebenen: operativ und strategisch}

In der Praxis trennt man zwei Tätigkeitsfelder. Die \emph{operative Beschaffung} kümmert
sich um die täglichen Routine-Aufgaben: Bestellungen auslösen, Lieferungen überwachen,
Reklamationen bearbeiten, Rechnungen prüfen. Hier zählt schnelle, fehlerfreie
Abwicklung. Die \emph{strategische Beschaffung} dagegen denkt langfristig: Welche
Lieferanten sollen Schlüsselpartner werden? Welche Materialien produzieren wir selbst,
welche kaufen wir ein? Wie verhandeln wir Rahmenverträge, die uns über Jahre günstige
Konditionen sichern? Beide Ebenen sind nötig --- ohne strategische Weichenstellung gerät
die operative Beschaffung in ständige Eilbestellungen, ohne saubere Operative wird die
beste Strategie nicht umgesetzt.

\subsection{Ziele der Beschaffung --- die „6~R”}

Damit das Tagesgeschäft nicht zur Beliebigkeit verkommt, formuliert die Betriebswirtschaft
sechs konkrete Ziele, an denen sich jede Beschaffungsentscheidung messen lassen muss.
Sie werden „die 6~R” genannt, weil jedes Ziel mit dem Wort \emph{richtig} beginnt.

\begin{merksatz}[title={Die sechs Beschaffungsziele}]
\textbf{Beschaffung soll das richtige Material, in der richtigen Menge, in der richtigen
Qualität, am richtigen Ort, zur richtigen Zeit, zum richtigen Preis bereitstellen.}
\end{merksatz}

\begin{eselsbruecke}
\textbf{„M-M-Q-O-Z-P”} --- \emph{Material, Menge, Qualität, Ort, Zeit, Preis}.
Wenn man die Anfangsbuchstaben in dieser Reihenfolge spricht, bleiben die sechs Ziele
leichter im Kopf.
\end{eselsbruecke}

Was bedeuten diese Ziele konkret? \textbf{Richtiges Material} heißt: das exakt benötigte
Produkt, eindeutig spezifiziert nach Art, Modell, Norm und Zustand. Wer „Aktenordner”
bestellt, statt „Aktenordner DIN~A4 breit, blau, mit Hebelmechanik”, riskiert, dass
das Falsche geliefert wird.

\textbf{Richtige Menge} bedeutet, weder zu viel noch zu wenig zu beschaffen. Zu wenig
führt zu Versorgungslücken und Stillstand, zu viel zu Lagerkosten und gebundenem
Kapital. Mengenrabatte locken zwar zu Großbestellungen, doch die Rechnung muss immer
\emph{Stückpreis-Ersparnis vs.\ Lagerkosten} enthalten.

\textbf{Richtige Qualität} ist nicht automatisch die höchste Qualität. Wer Premium-Foto­
papier für den Standard-Bürodruck einkauft, verschwendet Geld. Die richtige Qualität
ist diejenige, die der Verwendungszweck verlangt --- nicht mehr und nicht weniger.

\textbf{Richtiger Ort} verlangt Lieferung dorthin, wo das Material gebraucht wird ---
nicht ins Zentrallager, wenn die Filiale es nutzt. Moderne Just-in-Time-Konzepte
liefern bis direkt an die Produktionslinie.

\textbf{Richtige Zeit} heißt pünktlich. Zu früh verursacht Lagerkosten, zu spät führt
zu Produktionsstillstand oder verärgerten Kunden. Liefertermine gehören vertraglich
fixiert, gegebenenfalls mit Vertragsstrafen bei Verzug abgesichert.

\textbf{Richtiger Preis} ist nicht der billigste, sondern der wirtschaftlich beste.
Ein 5\,\% günstigeres Angebot mit 30\,\% Reklamationsquote ist am Ende teuer.

\subsection{Zielkonflikte}

Die 6~R stehen häufig in Spannung zueinander. Eine große Bestellmenge senkt den
Stückpreis (Ziel Preis), erhöht aber die Lagerkosten (Konflikt mit Wirtschaftlichkeit).
Ein billiger Lieferant (Ziel Preis) liefert vielleicht später (Konflikt mit Zeit) oder
schlechter (Konflikt mit Qualität). Eine Eilbestellung erfüllt das Ziel Zeit, aber nur
mit teurem Express­zuschlag (Konflikt mit Preis). Aus diesem Grund ist Beschaffung kein
einfaches Maximierungs-, sondern stets ein Optimierungs­problem: Welcher Mix der 6~R
ist im konkreten Fall am vorteilhaftesten?

\section{Organisation der Beschaffung}

\subsection{Zentral oder dezentral?}

Größere Unternehmen müssen entscheiden, ob sie ihre Beschaffung in einer zentralen
Einkaufsabteilung bündeln oder dezentral in den einzelnen Standorten oder Bereichen
ansiedeln.

\textbf{Zentrale Beschaffung} hat den Vorteil, dass sie große Mengen bündeln und damit
attraktive Mengenrabatte aushandeln kann. Sie schafft einheitliche Prozesse, bündelt
Einkaufs-Know-how an einer Stelle und ermöglicht eine konsequente Lieferantenstrategie.
Nachteil: Sie ist weiter vom konkreten Bedarf entfernt und reagiert oft langsamer.

\textbf{Dezentrale Beschaffung} hingegen sitzt nah am Geschehen. Die Filialleitung kann
schnell auf Engpässe reagieren, kennt die regionalen Lieferanten und Marktbedingungen.
Doch sie verschenkt Einkaufsmacht, weil viele kleine Bestellungen ungebündelt bleiben,
und neigt zu Wildwuchs in Standards und Konditionen.

\begin{tabular}{p{4.2cm} p{5.5cm} p{5.0cm}}
\toprule
\textbf{Kriterium} & \textbf{Zentrale Beschaffung} & \textbf{Dezentrale Beschaffung}\\
\midrule
Mengenrabatte & Hoch (Bündelung) & Niedrig\\
Marktnähe & Gering & Hoch\\
Spezialisierung & Einkaufsexperten & Verteiltes Wissen\\
Reaktionszeit & Langsam & Schnell vor Ort\\
Standardisierung & Einheitlich & Wildwuchs möglich\\
\bottomrule
\end{tabular}

\medskip
In der Praxis sieht man häufig einen \emph{Hybrid}: Die Zentrale verhandelt
Rahmenverträge mit Schlüssellieferanten, die einzelnen Bereiche rufen daraus dezentral
ab. Diese Mischform vereint die Vorteile beider Welten und ist heute Standard in
mittleren bis großen Unternehmen.

\subsection{Beschaffungswege}

Beim \textbf{direkten Bezug} kauft das Unternehmen unmittelbar beim Hersteller ein.
Das spart die Marge des Zwischenhändlers, ermöglicht direkte Qualitätsabsprachen und
ist meistens günstiger --- vorausgesetzt, die Bestellmenge ist groß genug, dass der
Hersteller überhaupt direkt verkauft.

Beim \textbf{indirekten Bezug} läuft die Bestellung über Großhandel, Vermittler oder
Importeur. Das ist teurer pro Stück, dafür kann man kleine Mengen abnehmen und hat
einen einzigen Ansprechpartner für ein breites Sortiment vieler Hersteller.

Im Büromanagement ist der indirekte Bezug für Standard-Verbrauchsmaterial die Regel:
Niemand bestellt 50 Aktenordner direkt beim Hersteller in Ostwestfalen, sondern beim
Bürofachgroßhandel um die Ecke oder beim Online-Anbieter wie Mercateo.

\section{Make or Buy --- Eigenfertigung oder Fremdbezug?}

Eine der grundlegendsten Entscheidungen der Beschaffung. Sie beantwortet die Frage,
\emph{ob} ein Sachgut oder eine Dienstleistung selbst erstellt oder von außen
hinzugekauft wird. Eine seriöse Make-or-Buy-Entscheidung kombiniert immer zwei
Perspektiven:

\begin{enumerate}
  \item \textbf{Quantitative Faktoren} --- Kostenvergleich. Welche Personal-, Material-,
    Verwaltungskosten fallen für die Eigenleistung an, im Vergleich zum Preis des
    Fremdbezugs?
  \item \textbf{Qualitative Faktoren} --- Strategie und Risiko. Verlieren wir Know-how?
    Werden wir vom Lieferanten abhängig? Können wir Qualität noch kontrollieren?
\end{enumerate}

\subsection{Beispielrechnung Reinigungsdienst}

Die fiktive \emph{Müller GmbH} überlegt, ihre Büroreinigung künftig durch einen externen
Dienstleister durchführen zu lassen. Heute hat sie eigene Reinigungskräfte beschäftigt.

\begin{center}
\begin{tabular}{p{6cm} r r}
\toprule
\textbf{Position} & \textbf{Make (€\,/\,Monat)} & \textbf{Buy (€\,/\,Monat)}\\
\midrule
Personalkosten Reinigung & 4\,200,00 & --- \\
Reinigungsmaterial & 350,00 & --- \\
Anteilige Verwaltungskosten & 280,00 & --- \\
Externer Dienstleister (Pauschalvertrag) & --- & 4\,500,00 \\
\midrule
\textbf{Gesamtkosten} & \textbf{4\,830,00} & \textbf{4\,500,00}\\
\midrule
\textbf{Differenz} & \multicolumn{2}{r}{\textbf{+\,330,00\,€\,/\,Monat zugunsten Buy ($\approx$\,6{,}8\,\%)}}\\
\bottomrule
\end{tabular}
\end{center}

Quantitativ ist der Fremdbezug also rund 330\,€ pro Monat günstiger, hochgerechnet
sind das knapp 4\,000\,€ im Jahr. Klingt nach klarer Empfehlung --- aber halt:
6,8\,\% Kostenvorteil ist sehr knapp. Bevor die Geschäftsführung wechselt, muss die
qualitative Bewertung folgen.

\subsection{Qualitative Bewertung des Beispiels}

Bei einem so geringen Kostenvorteil entscheiden die qualitativen Faktoren. Verliert
die Müller GmbH die direkte Qualitätskontrolle? Ist sie ab sofort vom Vertrag und der
Verfügbarkeit des Dienstleisters abhängig? Wer kümmert sich künftig um die
Reinigung-Sonderwünsche (Empfangsbereich vor Kundenbesuch besonders gründlich,
nach einer Veranstaltung am Wochenende)? Muss ein Vertragsmanagement aufgebaut
werden, das vorher nicht nötig war?

\begin{eselsbruecke}
\textbf{„Wenn der Vorteil klein ist, bleib's daheim.”} --- Bei knappem Kostenvorteil
im Fremdbezug überwiegen oft die qualitativen Risiken (Qualitätsverlust, Abhängigkeit,
Vertragsmanagement-Aufwand). Die endgültige Entscheidung muss beide Sichten
zusammenführen.
\end{eselsbruecke}

\subsection{Wann Make, wann Buy?}

Als Faustregel: Was \emph{strategisch wichtig} und \emph{nicht standardisierbar} ist,
sollte in der Regel selbst gemacht werden (Make). Was \emph{standardisiert} und
\emph{nicht differenzierend} ist, kann meistens günstiger fremdbezogen werden (Buy).
Reinigung ist nicht Kerngeschäft einer Industriefirma --- daher tendiert die Praxis
zum Buy. Forschung und Entwicklung dagegen bleibt fast immer im Haus.

\section{Beschaffungsprinzipien}

Welche Strategie bestimmt, \emph{wie} eingekauft wird? Drei Grundprinzipien stehen
zur Wahl --- in der Praxis kombiniert je Material und Risiko.

\subsection{Einzelbeschaffung im Bedarfsfall}

Hier wird genau dann eingekauft, wenn der konkrete Bedarf entsteht. Es gibt kein Lager,
keine Bevorratung. Der administrative Aufwand pro Bestellung ist hoch, aber dafür
fließt kein Kapital ins Lager. Klassisches Anwendungsfeld: Sondermaschinen,
Einzelfertigung, Sonder­anfertigungen --- alles, was nur einmal benötigt wird.

\subsection{Vorratsbeschaffung}

Hier werden größere Mengen auf Vorrat eingekauft und im Lager gehalten. Die
Versorgungssicherheit ist hoch, Mengenrabatte können genutzt werden, jedoch sind
Lagerkosten und Kapitalbindung deutlich. Im Büromanagement ist die Vorratsbeschaffung
der Normalfall für Standard-Verbrauchsmaterial wie Druckerpapier, Aktenordner oder
Schreibgeräte.

\subsection{Fertigungssynchrone Beschaffung (Just-in-Time)}

Just-in-Time geht einen Schritt weiter: Die Lieferung erfolgt exakt zum Bedarfszeit­
punkt --- weder zu früh noch zu spät. Lager im klassischen Sinn entfällt praktisch.
Vorteil: minimaler Bestand, keine Kapitalbindung. Nachteil: hohe Abhängigkeit von der
Lieferzuverlässigkeit. Ein einziger gestörter Lieferant kann die Produktion stoppen.

JIT funktioniert nur mit \emph{sehr zuverlässigen} Lieferanten, enger IT-Integration
(EDI-Anbindung, gemeinsame Bestandsdaten) und disziplinierten Prozessen. Klassisches
Anwendungsfeld: die Automobil­industrie. Toyota hat das Prinzip in den 1970er Jahren
zur Perfektion entwickelt.

Eine Variante ist das \textbf{Kanban}-Verfahren: Karten oder elektronische Signale
lösen die Nachbestellung automatisch aus, sobald ein Mindestbestand erreicht ist.

\section{E-Procurement}

\subsection{Was ist E-Procurement?}

\begin{merksatz}
E-Procurement (Electronic Procurement) ist die elektronische Unterstützung des
Beschaffungs­prozesses --- von der Bestellung bis zur Bezahlung. Es ersetzt papierbasierte
und telefonische Bestellungen durch IT-gestützte Workflows.
\end{merksatz}

Vorteile: schnellere Abwicklung, weniger Fehler, automatisierte Genehmigungs­
workflows, bessere Auswertbarkeit (Reporting). Nachteile: Investitions­kosten in
IT-Systeme, Abhängigkeit von der IT --- bei Ausfall steht der Bestellprozess.

\subsection{Drei Konzepte}

Drei Grundkonzepte unterscheiden sich danach, \emph{wer das System betreibt}:

\textbf{Sell-Side-Lösung} --- der \emph{Lieferant} betreibt einen Onlineshop, der
Käufer kauft dort ein. Beispiel: Ein Bürofachhandel betreibt einen Webshop, in dem
Firmenkunden ihre Bestellungen aufgeben. Vorteil für den Käufer: keine eigene IT
nötig. Nachteil: pro Lieferant ein eigenes System.

\textbf{Buy-Side-Lösung} --- der \emph{Käufer} betreibt ein eigenes Bestellsystem,
in dem zugelassene Lieferanten ihre Kataloge einstellen. Typisch bei großen Konzernen,
die ihre Beschaffung zentral steuern. Beispiel: SAP Ariba. Vorteil: volle Kontrolle
über Sortiment, Konditionen, Genehmigungs­prozesse.

\textbf{Marktplatz} --- ein neutraler Dritter bündelt Sortimente vieler Lieferanten.
Beispiel: Mercateo. Der Käufer hat eine breite Auswahl ohne eigene IT-Investition,
der Lieferant erreicht viele Käufer ohne eigenen Webshop.

\begin{eselsbruecke}
\textbf{„S~--~B~--~M”} --- \emph{Sell-Side} = Verkäufer-Webshop, \emph{Buy-Side} =
Käufer-Plattform, \emph{Marktplatz} = neutraler Bündler.
\end{eselsbruecke}

\section{Ökologische Beschaffung}

Nachhaltigkeit ist heute kein freiwilliger Bonus, sondern ein \emph{verpflichtendes
Beschaffungs­kriterium}. Zwei Treiber bestimmen die Entwicklung: gesetzliche Pflichten
und gesellschaftliche Erwartung.

\subsection{Rechtliche Treiber}

\textbf{Lieferketten­sorgfalts­pflichten­gesetz (LkSG)} --- seit 2023 verpflichtet das
Gesetz Unternehmen ab einer bestimmten Größe (zuletzt: ab 1\,000 Mitarbeitenden), in
ihrer Lieferkette Sorgfalt walten zu lassen: Menschenrechte achten, Umweltrisiken
identifizieren, dokumentieren, mindern. Verstöße können zu hohen Bußgeldern führen.

\textbf{EU-Lieferketten­richtlinie (CSDDD)} --- erweitert diese Sorgfaltspflichten
europaweit. In Deutschland wird sie 2026/2027 in nationales Recht umgesetzt und
verschärft die Anforderungen.

\textbf{Corporate Sustainability Reporting Directive (CSRD)} --- verpflichtet
Unternehmen zur transparenten Berichterstattung über CO\textsubscript{2}-Bilanz,
Soziales und Governance entlang der Lieferkette. Wer nicht messen kann, kann nicht
berichten.

\subsection{Bewertungs­kriterien für nachhaltige Lieferanten}

\begin{itemize}[itemsep=0pt]
  \item \textbf{Transportwege:} kurze Lieferwege, regionale Bezugsquellen
  \item \textbf{Verpackung:} reduziert, recyclingfähig, mehrweg
  \item \textbf{Energie­bilanz:} Anteil erneuerbarer Energien beim Lieferanten
  \item \textbf{CO\textsubscript{2}-Fußabdruck:} pro Produkt, pro Lieferung
  \item \textbf{Sozial­standards:} Faire Löhne, sichere Arbeitsbedingungen
  \item \textbf{Recyclingfähigkeit:} Kreislauffähigkeit am Lebenszyklusende
\end{itemize}

In der Praxis prüfen Käufer ihre Lieferanten zunehmend in \emph{Lieferanten-Audits}
(Vor-Ort-Prüfungen) und verlangen Zertifikate (z.\,B.\ ISO 14001 für Umweltmanagement,
SA 8000 für Sozialstandards).

\section{Zusammenfassung}

\begin{enumerate}
  \item Beschaffung ist Bindeglied zwischen Markt und Produktion in der Wertschöpfungskette.
  \item Die \textbf{6~R} (Material, Menge, Qualität, Ort, Zeit, Preis) beschreiben die operativen Ziele und stehen häufig im Konflikt.
  \item \textbf{Make or Buy} kombiniert Kostenvergleich (quantitativ) mit strategischen Faktoren (qualitativ).
  \item Drei \textbf{Beschaffungsprinzipien}: Einzel, Vorrats, fertigungssynchron --- je nach Material und Strategie.
  \item \textbf{E-Procurement} in drei Varianten: Sell-Side, Buy-Side, Marktplatz.
  \item \textbf{Ökologische Beschaffung} ist heute Pflicht (LkSG, CSRD), nicht Kür.
  \item Der Prozess folgt 7 Schritten: Bedarf $\to$ Quelle $\to$ Anfrage $\to$ Vergleich $\to$ Bestellung $\to$ Wareneingang $\to$ Rechnung\,\&\,Zahlung.
\end{enumerate}

\newpage

\section*{Glossar}
\addcontentsline{toc}{section}{Glossar}

\glossareintrag{Beschaffung}{Alle Tätigkeiten, durch die ein Unternehmen die zur Leistungserstellung benötigten Sachgüter und Dienstleistungen bereitstellt.}

\glossareintrag{Wertschöpfungskette}{Folge von Tätigkeiten, die einen Wertzuwachs vom Material bis zum verkaufsfertigen Produkt bewirken.}

\glossareintrag{6~R}{Die sechs operativen Beschaffungsziele: richtiges Material, richtige Menge, richtige Qualität, richtiger Ort, richtige Zeit, richtiger Preis.}

\glossareintrag{Operative Beschaffung}{Tagesgeschäft: Bestellungen, Lieferüberwachung, Reklamationen, Rechnungsprüfung.}

\glossareintrag{Strategische Beschaffung}{Langfristige Aufgaben: Lieferanten-, Material-, Make-or-Buy-Strategie, Rahmenverträge.}

\glossareintrag{Make or Buy}{Strategische Entscheidung zwischen Eigenfertigung (Make) und Fremdbezug (Buy) eines Sachgutes oder einer Dienstleistung.}

\glossareintrag{Einzelbeschaffung}{Beschaffungsprinzip, bei dem pro Auftrag gezielt eingekauft wird; geringer Lagerbestand, hoher administrativer Aufwand.}

\glossareintrag{Vorratsbeschaffung}{Beschaffungsprinzip, bei dem größere Mengen auf Vorrat eingekauft werden; hohe Versorgungssicherheit, hohe Kapitalbindung.}

\glossareintrag{Just-in-Time (JIT)}{Fertigungssynchrones Beschaffungsprinzip; Material wird genau zum Bedarfszeitpunkt geliefert, Lager praktisch null.}

\glossareintrag{Kanban}{Aus Japan stammendes JIT-Verfahren mit Karten/Signalen, das den Materialfluss bedarfsgesteuert auslöst.}

\glossareintrag{E-Procurement}{Elektronische Unterstützung des Beschaffungsprozesses; drei Konzepte: Sell-Side, Buy-Side, Marktplatz.}

\glossareintrag{Sell-Side-Lösung}{E-Procurement-System, das vom \emph{Lieferanten} betrieben wird (z.\,B.\ Onlineshop).}

\glossareintrag{Buy-Side-Lösung}{E-Procurement-System, das vom \emph{Käufer} betrieben wird; Lieferanten stellen Kataloge ein.}

\glossareintrag{Marktplatz}{E-Procurement-System eines neutralen Dritten, der Sortimente vieler Lieferanten bündelt.}

\glossareintrag{LkSG}{Lieferketten\-sorgfalts\-pflichten\-gesetz; verpflichtet ab 2023 deutsche Unternehmen zur Sorgfalt entlang ihrer Lieferkette.}

\glossareintrag{CSDDD}{Corporate Sustainability Due Diligence Directive --- EU-Lieferkettenrichtlinie, die LkSG europaweit erweitert.}

\glossareintrag{CSRD}{Corporate Sustainability Reporting Directive; EU-Richtlinie zur verpflichtenden Nachhaltigkeits\-bericht\-erstattung.}

\glossareintrag{Mengenrabatt}{Preisnachlass bei Abnahme größerer Mengen; ein Anreiz zur Vorratsbeschaffung.}

\glossareintrag{Skaleneffekte}{Stückkostenvorteil bei größeren Produktionsmengen; ein Argument für Fremdbezug bei Spezialisten.}

\glossareintrag{Lieferanten-Audit}{Formale Prüfung eines Lieferanten auf Einhaltung definierter Standards (Qualität, Umwelt, Soziales).}

\glossareintrag{Rahmenvertrag}{Langfristiger Vertrag, der Konditionen für eine Reihe von Einzelabrufen festlegt.}

\end{document}
